\newpage
\Solution{4}
If $\vv{F}_1$ and $\vv{F}_2$ are parallel, then there exists some vector $\vv{u}$ which is perpendicular to both:
\begin{equation*}
	\vv{u} \cdot \vv{F}_1 = 0
	\quad\text{and}\quad
	\vv{u} \cdot \vv{F}_2 = 0
\end{equation*}
Then for the resultant force we show:
\begin{equation*}
	\vv{u} \cdot \vv{F}_{res}
	=
	\vv{u} \cdot 
	\left(
	\vv{F}_1 + \vv{F}_2
	\right)
	=
	\vv{u} \cdot \vv{F}_1
	+
	\vv{u} \cdot \vv{F}_2
	=
	0
	\quad\Rightarrow\quad
	\boxed{\vv{u} \cdot \vv{F}_{res} = 0}
	\tag{$1$}
	\label{eq:j17_p4_1}
\end{equation*}
Since $\vv{F}_{res}$ is also perpendicular to $\vv{u}$, then we have shown that it is parallel to both $\vv{F}_1$ and $\vv{F}_2$. Given this, the vector diagram for $\vv{F}_{res} = \vv{F}_1 + \vv{F}_2$ degenerates to a line segment. Therefore:
\begin{equation*}
	\boxed{F_{res} = F_1 + F_2
	}
	\tag{$2$}
	\label{eq:j17_p4_2}
\end{equation*}
\begin{minipage}[t][145pt]{\textwidth}
	\vspace{0pt}
	\begin{minipage}[t][143pt]{0.5\textwidth}
	\vspace{0pt}
	\centering
	\begin{tikzpicture}
		% Forces
		\node[point] at (3, 6) {};
		\node[point] at (6, 10) {};
		
		\node[below left] at (3, 6) {$A$};
		\node[above] at (6, 10) {$B$};
		
		\draw[line, -{Stealth}] (3, 6) -- (5.5, 6);
		\draw[line, -{Stealth}] (6, 10) -- (7.5, 10);
		
		\node[right] at (5.5, 6) {$\vv{F_1}$};
		\node[right] at (7.5, 10) {$\vv{F_2}$};
		
		% Resultant force
		\draw[line, -{Stealth}] (5.25, 9) -- (9, 9);
		\node[point] at (5.25, 9) {};
		\node[right] at (9, 9) {$\vv{F}_{res}$};
		
		% CM line
		\draw[line, dashed] (3, 7.5) -- (8, 7.5);
		
		\node[above] at (7, 7.5) {\text{CM line of}};
		\node[below] at (7, 7.5) {\text{action}};
		
		% Triangle
		\draw[line] (3, 6) -- (6, 10);
		\draw[line, dashed] (3, 6) -- (3, 10);
		\draw[line, dashed] (3, 10) -- (6, 10);
		\draw[line, dashed] (5.25, 9) -- (3, 9);
		
		\node[point] at (4.125, 7.5) {};
		\node[point] at (3, 7.5) {};
		\node[point] at (3, 10) {};
		\node[point] at (3, 9) {};
		
		\node[below] at (4.2, 7.5) {$C$};
		\node[left] at (3, 7.5) {$M$};
		\node[left] at (3, 10) {$N$};
		\node[above right] at (3, 9) {$T$};
		\node[above] at (5.2, 9) {$S$};
		
		% Brackets and dimensions
		\draw[bracket] (3, 6) -- (3, 7.5)
		node[midway, left=5pt] {$x$};
		\draw[bracket] (3, 7.5) -- (3, 10)
		node[midway, left=5pt] {$y$};
		\draw[line, Stealth-Stealth] (3.5, 7.55) -- (3.5, 8.95)
		node[midway, right] {$c$};
	\end{tikzpicture}
\end{minipage}
	\begin{minipage}[t][143pt]{0.49\textwidth}
		\vspace{2pt}
		For the last part consider the shown diagram. Assume that the distances from the lines of action of $\vv{F}_1$ and $\vv{F}_2$ to the line of action of the centre of mass are $x$ and $y$ respectively. In addition, let the distance from the line of action of $\vv{F}_{res}$ to the centre of mass line of action be $c$. From the torque law:
		\begin{equation*}
			F_2 y - F_1 x = F_{res} c = (F_1 + F_2) c
		\end{equation*}
		Some algebra yields:
		\begin{equation*}
			F_1 (x + c) = F_2 (y - c)
		\end{equation*}
	\end{minipage}
\end{minipage}
Notice that $x + c = AT$ and $y - c = TN$. Then we have:
\begin{equation*}
	\boxed{\frac{AT}{TN} = \frac{F_2}{F_1} = \frac{AS}{SB}}
	\tag{$3$}
	\label{eq:j17_p4_3}
\end{equation*}
where the last equality follows from Thales' Theorem.